\begin{table*}[htbp]
\centering
\caption{Geophysical context of the real-observation reconstruction cases. Solar-wind and IMF values are the conditioning values used for each case; AE and SYM-H summarize the corresponding context window. Kp is obtained from the 3-hour OMNI2 index when available.}
\label{tab:event_context}
\scriptsize
\resizebox{\textwidth}{!}{%
\begin{tabular}{lllllrrrrrrrr}
\hline
Event & Instrument & Time / interval (UT) & Orbit & Polar-region crossing (UT) & $B_x$ & $B_y$ & $B_z$ & $V_{\mathrm{sw}}$ & $P_{\mathrm{dyn}}$ & Kp & AE$_{\max}$ & SYM-H$_{\min}$ \\
 & & & & & \multicolumn{3}{c}{(nT)} & (km s$^{-1}$) & (nPa) & & (nT) & (nT) \\
\hline
Figure 4 & Polar/UVI & 1996-04-03 02:15 & 0 & 20:15--04:38 & 0.9 & -1.0 & -1.6 & 325 & 2.91 & 1.3 & 23 & 1 \\
Figure 5 & Polar/UVI & 1996-04-01 08:40 & 0 & 08:42--14:30 & -2.9 & -3.7 & -3.8 & 365 & 2.69 & 1.7 & 337 & -18 \\
Figure 6 & DMSP F16/SSUSI & 2005-01-01 14:44:34--16:26:25 & 06228 & 14:44--16:26 & -6.5 & 7.7 & 1.7 & 459 & 3.18 & 3.3 & 150 & 8 \\
Figure 7 & DMSP F16/SSUSI & 2005-01-04 02:12:15--03:54:07 & 06263 & 02:12--03:54 & -4.5 & -2.0 & -2.1 & 656 & 2.89 & 4.3 & 515 & -28 \\
\hline
\end{tabular}%
}
\end{table*}
